\documentclass{labreport}

\usepackage{microtype}
\usepackage{enumitem}
\usepackage{pifont}
\usepackage{tabularx}
\usepackage{booktabs}

% Colors
\definecolor{codebg}{RGB}{248, 249, 251}
\definecolor{codeborder}{RGB}{218, 222, 229}
\definecolor{linkcolor}{RGB}{20, 60, 120}
\definecolor{zshgreen}{RGB}{0, 160, 80}
\definecolor{zshred}{RGB}{200, 30, 30}

% Hyperref setup with subtle link colors
\hypersetup{
    colorlinks=true,
    linkcolor=linkcolor,
    filecolor=linkcolor,
    urlcolor=linkcolor,
    citecolor=linkcolor
}

% Code and terminal listings
\lstdefinestyle{bash}{
    basicstyle=\ttfamily\footnotesize,
    backgroundcolor=\color{codebg},
    frame=single,
    framerule=0.5pt,
    rulecolor=\color{codeborder},
    breaklines=true,
    breakatwhitespace=false,
    showstringspaces=false,
    keepspaces=true,
    columns=flexible,
    numbers=none,
    xleftmargin=3pt,
    xrightmargin=3pt,
    framesep=4pt
}

\lstdefinestyle{c}{
    language=C,
    basicstyle=\ttfamily\footnotesize,
    backgroundcolor=\color{codebg},
    frame=single,
    framerule=0.5pt,
    rulecolor=\color{codeborder},
    keywordstyle=\color{blue!85!black}\bfseries,
    commentstyle=\color{gray!80!black}\itshape,
    stringstyle=\color{teal!80!black},
    numbers=left,
    numberstyle=\tiny\color{gray!80!black},
    numbersep=8pt,
    breaklines=true,
    breakatwhitespace=false,
    showstringspaces=false,
    keepspaces=true,
    columns=flexible,
    tabsize=2,
    xleftmargin=14pt,
    xrightmargin=3pt,
    framesep=4pt
}

\lstset{
    extendedchars=true,
    literate=
        {➜}{{\textcolor{zshgreen}{\ding{220}}}}{1}
        {├─}{|-\,}{2}
        {└─}{\`-\,}{2}
        {├}{|}{1}
        {─}{-}{1}
        {└}{\`}{1}
        {✗}{{\textcolor{zshred}{\ding{55}}}}{1}
}

% Section styling for clean document structure
\titleformat{\section}
  {\large\bfseries}{\thesection}{0.8em}{}[\vspace{1pt}\hrule height 0.5pt]
\titleformat{\subsection}
  {\normalsize\bfseries}{\thesubsection}{0.8em}{}
\titleformat{\subsubsection}
  {\small\bfseries}{\thesubsubsection}{0.8em}{}

\begin{document}

\labtitle
    {4}
    {The xv6 Shell from the Outside}
    {st. gr. FAF-243}
    {Șoimu Ionuț}
    {asist. univ.\\[0.15cm] Reițman Patricia}

\tableofcontents
\newpage

% ==============================================================================
\section{Lab 4: Track B -- The xv6 Shell from the Outside}

\subsection{Before you start}

\subsubsection*{Purpose and Scope}
The entire user interface of the educational UNIX operating system xv6 is implemented in a single user-space program, \texttt{user/sh.c}, spanning approximately 500 lines of concise C code. The primary purpose of this laboratory work is to investigate the mechanics of command-line interpretation by running the xv6 shell side-by-side with modern GNU \texttt{bash} on Linux.

Through experimental side-by-side comparison, we dissect what capabilities an operating system shell is fundamentally required to provide (process invocation, file descriptor redirection, and inter-process pipelining) versus what conveniences modern shells like \texttt{bash} layer on top (pathname globbing, environment variable expansion, quotation parsing, and interactive line-editing). Furthermore, by writing a dedicated argument-inspection utility (\texttt{args}) in both environments, we empirically prove where the boundary of responsibility lies between user-space utilities, the shell interpreter, and the kernel.

\subsubsection*{Environment Setup and Identification}
\begin{lstlisting}[style=bash]
➜  ~ cd ~/xv6-riscv
➜  xv6-riscv git:(riscv) whoami
ionut
➜  xv6-riscv git:(riscv) uname -a
Linux arch 7.2.8-1-cachyos #1 SMP PREEMPT_DYNAMIC Sat, 26 Sep 2026 19:43:45 +0000 x86_64 GNU/Linux
➜  xv6-riscv git:(riscv) git status
On branch riscv
Your branch is up to date with 'origin/riscv'.
\end{lstlisting}

% ------------------------------------------------------------------------------
\subsection{Part 1. Fix sleep for the current xv6}

\subsubsection{System Call Verification in user.h}
In August 2025, the xv6 upstream development team renamed the timer system call \texttt{sleep()} to \texttt{pause()} to avoid conceptual naming collisions with POSIX sleep semantics. We inspect \texttt{user/user.h} to determine which system call declaration is present:

\begin{lstlisting}[style=bash]
➜  xv6-riscv git:(riscv) grep -n "pause\|sleep" user/user.h
25:int pause(int);
\end{lstlisting}
The header file explicitly declares \texttt{int pause(int);}. Consequently, our user-space utility must invoke \texttt{pause()} to enter the sleep state.

\subsubsection{Implementation of user/sleep.c}
We implement \texttt{user/sleep.c}, ensuring proper argument validation and usage diagnostics:

\begin{lstlisting}[style=c]
#include "kernel/types.h"
#include "user/user.h"

int
main(int argc, char *argv[])
{
  if(argc != 2){
    fprintf(2, "usage: sleep <ticks>\n");
    exit(1);
  }
  pause(atoi(argv[1]));   // system call into the kernel (was sleep() before Aug 2025)
  exit(0);
}
\end{lstlisting}

\subsubsection{Integration into Makefile}
We ensure that the binary target \texttt{\$U/\_sleep\textbackslash} is registered within the \texttt{UPROGS} definition list in the xv6 build configuration:
\begin{lstlisting}[style=bash]
UPROGS=\
	$U/_cat\
	$U/_echo\
	$U/_sleep\
	$U/_forktest\
	...
\end{lstlisting}

\subsubsection{Build and Testing in QEMU}
We compile the operating system and user-space binaries and boot the kernel inside QEMU:
\begin{lstlisting}[style=bash]
➜  xv6-riscv git:(riscv) make qemu
qemu-system-riscv64 -machine virt -bios none -kernel kernel/kernel -m 128M -smp 3 -nographic -global virtio-mmio.force-legacy=false -drive file=fs.img,if=none,format=raw,id=x0 -device virtio-blk-device,drive=x0,bus=virtio-mmio-bus.0

xv6 kernel is booting

hart 1 starting
hart 2 starting
init: starting sh
$ sleep 20
$ sleep
usage: sleep <ticks>
$
\end{lstlisting}
The tests succeed: \texttt{sleep 20} pauses the shell for 20 timer ticks before re-prompting, and invoking \texttt{sleep} without arguments triggers the error message on file descriptor 2 with exit code 1.

% ------------------------------------------------------------------------------
\subsection{Part 2. bash versus the xv6 shell}

\subsubsection{Comparative Test Suite}
We execute a standardized suite of 10 commands in both shells:
\begin{itemize}[noitemsep]
    \item \textbf{Linux bash:} Executed inside an isolated directory \texttt{\~/shtest} containing an empty file named \texttt{README}.
    \item \textbf{xv6 shell:} Executed at the xv6 interactive prompt \texttt{\$} inside QEMU.
\end{itemize}

\newpage
\begin{table}[htbp]
\centering
\footnotesize
\renewcommand{\arraystretch}{1.12}
\setlength{\tabcolsep}{4pt}
\begin{tabularx}{\textwidth}{p{3.4cm} p{2.0cm} p{2.5cm} p{2.5cm} X}
\toprule
\textbf{Command} & \textbf{Tests} & \textbf{bash (Linux)} & \textbf{xv6 shell} & \textbf{Why / Explanation} \\
\midrule
\texttt{echo hello > f ; cat < f} & Redirection & \texttt{hello} & \texttt{hello} & Both shells redirect stdout (\texttt{>}) and stdin (\texttt{<}) using OS file descriptors. \\
\addlinespace
\texttt{ls | grep READ} & Pipeline & \texttt{README} & \texttt{README 2 2 2441} & Both support pipes (\texttt{|}). xv6 \texttt{ls} includes type, inode, and byte size. \\
\addlinespace
\texttt{(echo a; echo b) | wc} & Grouping, list & \texttt{2 2 4} & \texttt{2 2 4} & Both shells parse parentheses (\texttt{(...)}) as command blocks whose combined stream is piped. \\
\addlinespace
\texttt{echo one ; echo two} & Sequence & \texttt{one} \newline \texttt{two} & \texttt{one} \newline \texttt{two} & Both shells use \texttt{;} to execute commands sequentially. \\
\addlinespace
\texttt{echo bg \&} & Background & \texttt{bg} (asynchronous) & \texttt{bg} (asynchronous) & Both support \texttt{\&} by forking without calling \texttt{wait()} in the parent. \\
\addlinespace
\texttt{ls *} & Globbing & \texttt{args args\_linux.c f README} & \texttt{ls: cannot open *} & bash expands \texttt{*} to matching filenames; xv6 passes \texttt{*} literally. \\
\addlinespace
\texttt{echo \$HOME} & Variables & \texttt{/home/ionut} & \texttt{\$HOME} & bash expands environment variables; xv6 has no variables and prints \texttt{\$HOME} literally. \\
\addlinespace
\texttt{echo "a b"} & Quotes & \texttt{a b} & \texttt{"a b"} & bash removes quotes and groups words; xv6 treats quotes as normal characters. \\
\addlinespace
\texttt{mkdir d ; cd d ; ls ..} & PATH lookup & \texttt{args ... README} & \texttt{exec ls failed} & bash resolves \texttt{ls} via \texttt{\$PATH} (\texttt{/usr/bin}); xv6 only checks current directory \texttt{/d}. \\
\addlinespace
\texttt{Arrow Up, Tab, Ctrl+R} & History / Completion & Recalls command, autocompletes, searches & Prints escape codes / tabs; no history & bash links against GNU Readline; xv6 reads raw characters via \texttt{gets()}. \\
\bottomrule
\end{tabularx}
\caption{Comprehensive behavioral comparison between GNU bash and xv6 sh}
\end{table}

\subsubsection{Terminal Execution Record: GNU bash}
\begin{lstlisting}[style=bash]
➜  ~ mkdir -p ~/shtest && cd ~/shtest && touch README
➜  shtest echo hello > f ; cat < f
hello
➜  shtest ls | grep READ
README
➜  shtest (echo a; echo b) | wc
      2       2       4
➜  shtest echo one ; echo two
one
two
➜  shtest echo bg &
[1] 23412
bg
➜  shtest ls *
args  args_linux.c  f  README
➜  shtest echo $HOME
/home/ionut
➜  shtest echo "a b"
a b
➜  shtest mkdir d ; cd d ; ls ..
args  args_linux.c  d  f  README
➜  shtest cd ~/shtest
\end{lstlisting}

\subsubsection{Terminal Execution Record: xv6 shell}
\begin{lstlisting}[style=bash]
$ echo hello > f ; cat < f
hello
$ ls | grep READ
README         2 2 2441
$ (echo a; echo b) | wc
2 2 4
$ echo one ; echo two
one
two
$ echo bg &
bg
$ ls *
ls: cannot open *
$ echo $HOME
$HOME
$ echo "a b"
"a b"
$ mkdir d ; cd d ; ls ..
exec ls failed
$ /ls ..
.              1 1 1024
..             1 1 1024
README         2 2 2441
cat            2 3 35864
echo           2 4 34760
...
d              1 24 32
f              2 25 6
$
\end{lstlisting}

\subsubsection*{Observe:}
\begin{enumerate}[leftmargin=*]
    \item \textbf{List the features both shells have, and the ones only bash has.}
    \begin{itemize}[noitemsep, topsep=2pt]
        \item \textbf{Features common to both shells:}
        \begin{itemize}[noitemsep]
            \item Input, output, and appending redirection (\texttt{<}, \texttt{>}, \texttt{>>})
            \item Inter-process communication pipelines (\texttt{|})
            \item Command sequencing via semicolons (\texttt{;})
            \item Command grouping in subshells (\texttt{(...)})
            \item Asynchronous background process execution (\texttt{\&})
            \item Built-in directory navigation (\texttt{cd})
        \end{itemize}
        \item \textbf{Features present only in bash:}
        \begin{itemize}[noitemsep]
            \item \texttt{\$PATH} binary resolution across system directories
            \item Pathname expansion and wildcard globbing (\texttt{*}, \texttt{?}, \texttt{[...]})
            \item Environment and shell variables (\texttt{\$HOME}, \texttt{\$USER}, \texttt{\$?})
            \item Quote parsing, parameter grouping, and escape sequences (\texttt{""}, \texttt{''}, \texttt{\textbackslash})
            \item Command-line editing, history browsing, reverse search, and Tab completion (GNU Readline)
            \item Control-flow statements (\texttt{if}, \texttt{for}, \texttt{while}) and command substitution (\texttt{\$(...)})
        \end{itemize}
    \end{itemize}

    \item \textbf{In xv6, ls stops working after cd d. Why? Find a way to run ls from inside d anyway.}\\[0.2em]
    \textbf{Why it fails:} The xv6 shell does not maintain or search a \texttt{PATH} environment variable. When a command such as \texttt{ls} is executed without a leading slash, the shell directly invokes the system call \texttt{exec("ls", ...)}. The xv6 kernel attempts to open and load the executable named \texttt{ls} relative to the process's current working directory. When the user is in the root directory \texttt{/}, \texttt{ls} succeeds because the binary is located at \texttt{/ls}. However, once inside directory \texttt{/d}, the kernel searches for \texttt{/d/ls}, which does not exist, causing \texttt{exec} to fail with \texttt{exec ls failed}.\\
    \textbf{Way to run it anyway:} We must supply the absolute path or relative path to the root binary:
    \begin{center}
    \texttt{/ls ..} \quad \text{or} \quad \texttt{../ls ..}
    \end{center}

    \item \textbf{Which missing feature would you miss most in daily work, and which would be easiest to add?}\\[0.2em]
    \textbf{Missed most:} The lack of \textbf{\texttt{PATH} lookup} and \textbf{command-line editing (Tab completion / history)}. Without \texttt{PATH}, organizing programs into clean directories like \texttt{/bin} or \texttt{/usr/bin} is impractical without typing full paths for every utility. Without Tab completion and history, every path and typo must be re-entered manually from scratch.\\
    \textbf{Easiest to add:} A basic \textbf{\texttt{PATH} lookup mechanism} or \textbf{simple quote stripping}. For \texttt{PATH}, in \texttt{user/sh.c}, inside \texttt{case EXEC:}, if \texttt{ecmd->argv[0]} does not begin with \texttt{/}, the shell could simply prefix \texttt{/} to the program name before invoking \texttt{exec()}. If the file exists in the root directory, it executes immediately without requiring the user to type \texttt{/ls}.
\end{enumerate}

% ------------------------------------------------------------------------------
\subsection{Part 3. Who expands *? A program that prints its argv}

\subsubsection{Utility Implementation for xv6: user/args.c}
We write \texttt{user/args.c} to print the exact number of command-line arguments received and each individual string in \texttt{argv}:

\begin{lstlisting}[style=c]
#include "kernel/types.h"
#include "user/user.h"

int
main(int argc, char *argv[])
{
  printf("argc = %d\n", argc);
  for(int i = 0; i < argc; i++)
    printf("argv[%d] = [%s]\n", i, argv[i]);
  exit(0);
}
\end{lstlisting}
We add \texttt{\$U/\_args\textbackslash} to the \texttt{UPROGS} list in \texttt{Makefile}.

\subsubsection{Utility Implementation for Linux: args\_linux.c}
We write the equivalent C utility for Linux adhering to POSIX standards:

\begin{lstlisting}[style=c]
#include <stdio.h>

int
main(int argc, char *argv[])
{
  printf("argc = %d\n", argc);
  for(int i = 0; i < argc; i++)
    printf("argv[%d] = [%s]\n", i, argv[i]);
  return 0;
}
\end{lstlisting}
We compile the program using \texttt{gcc}:
\begin{lstlisting}[style=bash]
➜  shtest gcc -o args args_linux.c
\end{lstlisting}

\subsubsection{Side-by-Side Execution}
We run the exact same command lines in both shells:

\begin{lstlisting}[style=bash]
# ------------------------------------------------------------- GNU bash (Linux)
➜  shtest ./args one    two    three
argc = 4
argv[0] = [./args]
argv[1] = [one]
argv[2] = [two]
argv[3] = [three]

➜  shtest ./args *.c "x y"
argc = 3
argv[0] = [./args]
argv[1] = [args_linux.c]
argv[2] = [x y]

➜  shtest ./args $HOME
argc = 2
argv[0] = [./args]
argv[1] = [/home/ionut]

➜  shtest ./args 'single quoted'
argc = 2
argv[0] = [./args]
argv[1] = [single quoted]
\end{lstlisting}

\begin{lstlisting}[style=bash]
# ------------------------------------------------------------- xv6 shell (QEMU)
$ args one    two    three
argc = 4
argv[0] = [args]
argv[1] = [one]
argv[2] = [two]
argv[3] = [three]

$ args *.c "x y"
argc = 4
argv[0] = [args]
argv[1] = [*.c]
argv[2] = ["x]
argv[3] = [y"]

$ args $HOME
argc = 2
argv[0] = [args]
argv[1] = [$HOME]

$ args 'single quoted'
argc = 3
argv[0] = [args]
argv[1] = ['single]
argv[2] = [quoted']
\end{lstlisting}

\subsubsection*{Observe:}
\begin{enumerate}[leftmargin=*]
    \item \textbf{Paste both outputs for args *.c "x y". How many arguments did each program receive?}
    \begin{itemize}[noitemsep, topsep=2pt]
        \item \textbf{GNU bash output:}
\begin{lstlisting}[style=bash]
argc = 3
argv[0] = [./args]
argv[1] = [args_linux.c]
argv[2] = [x y]
\end{lstlisting}
        The Linux program received \textbf{3 arguments} (\texttt{argc = 3}).
        \item \textbf{xv6 shell output:}
\begin{lstlisting}[style=bash]
argc = 4
argv[0] = [args]
argv[1] = [*.c]
argv[2] = ["x]
argv[3] = [y"]
\end{lstlisting}
        The xv6 program received \textbf{4 arguments} (\texttt{argc = 4}).
    \end{itemize}

    \item \textbf{So who turns *.c into file names and removes the quotes: the program, the kernel, or the shell?}\\[0.2em]
    \textbf{The shell.} This experiment provides conclusive proof.\\
    Neither the target binary (\texttt{args}) nor the kernel performs wildcard expansion or quotation stripping. In both Linux and xv6, the binary code of \texttt{args} does nothing more than iterate over whatever pointers are already populated in \texttt{argv}. The kernel's \texttt{exec} system call simply transfers the array of string pointers given to it into the newly loaded process's address space without altering their contents.\\
    In Linux, \texttt{bash} inspects the line before invoking \texttt{exec}, expands \texttt{*.c} against filesystem entries, strips the enclosing double quotes from \texttt{"x y"}, and treats \texttt{x y} as a single argument containing a space.\\
    In xv6, \texttt{user/sh.c} does not possess globbing or quotation-parsing logic; it merely tokenizes the line on whitespace. Therefore, \texttt{*.c}, \texttt{"x}, and \texttt{y"} are handed directly to the kernel and onto \texttt{args} as raw, unmodified strings.

    \item \textbf{What happened to the extra spaces in one    two    three?}\\[0.2em]
    In both shells, consecutive spaces are treated strictly as \textbf{token delimiters}. When the shell's lexical scanner encounters whitespace, it consumes all contiguous spaces until the next non-whitespace character is reached. The tokens (\texttt{one}, \texttt{two}, and \texttt{three}) are extracted and stored into distinct array slots in \texttt{argv}; the whitespace characters themselves are discarded and never passed to the invoked program.
\end{enumerate}

% ------------------------------------------------------------------------------
\subsection{Part 4. Read the shell}

\subsubsection{Source Code Inspection of user/sh.c}
Inside the xv6 repository (\texttt{\~/xv6-riscv}), we inspect the structure and implementation details of the shell:

\begin{lstlisting}[style=bash]
➜  xv6-riscv git:(riscv) wc -l user/sh.c
499 user/sh.c

➜  xv6-riscv git:(riscv) grep -n "#define" user/sh.c
8:#define EXEC  1
9:#define REDIR 2
10:#define PIPE  3
11:#define LIST  4
12:#define BACK  5
14:#define MAXARGS 10

➜  xv6-riscv git:(riscv) grep -n "case " user/sh.c
75:  case EXEC:
83:  case REDIR:
93:  case LIST:
101:  case PIPE:
125:  case BACK:
284:  case 0:
286:  case '|':
287:  case '(':
288:  case ')':
289:  case ';':
290:  case '&':
291:  case '<':
294:  case '>':
391:    case '<':
394:    case '>':
397:    case '+': // >>
469:  case EXEC:
475:  case REDIR:
481:  case PIPE:
487:  case LIST:
493:  case BACK:

➜  xv6-riscv git:(riscv) grep -n "cd" user/sh.c | head -5
166:    if (cmd[0] == 'c' && cmd[1] == 'd' && cmd[2] == ' ') {
167:      // Chdir must be called by the parent, not the child.
168:      cmd[strlen(cmd) - 1] = 0; // chop \n
169:      if (chdir(cmd + 3) < 0)
170:        fprintf(2, "cannot cd %s\n", cmd + 3);
\end{lstlisting}

\subsubsection{Architecture of the xv6 Command Interpreter}
The parser in \texttt{user/sh.c} models commands as an abstract syntax tree of polymorphic command structures sharing a common header (\texttt{struct cmd \{ int type; \}}):
\begin{itemize}[noitemsep]
    \item \texttt{struct execcmd}: Holds program arguments (\texttt{argv}) for \texttt{exec()}.
    \item \texttt{struct redircmd}: Holds target file, flags, and file descriptor for redirection.
    \item \texttt{struct pipecmd}: Contains \texttt{left} and \texttt{right} subcommands connected via an anonymous pipe.
    \item \texttt{struct listcmd}: Contains subcommands executed in sequence (\texttt{;}).
    \item \texttt{struct backcmd}: Wraps a subcommand to be spawned asynchronously (\texttt{\&}).
\end{itemize}

\subsubsection*{Observe:}
\begin{enumerate}[leftmargin=*]
    \item \textbf{Name the five kinds of command the xv6 shell knows and the character that produces each.}
    \begin{enumerate}[label=\arabic*., noitemsep, topsep=2pt]
        \item \textbf{\texttt{EXEC}} (Type 1): An ordinary executable command; produced by regular word tokens delimited by whitespace.
        \item \textbf{\texttt{REDIR}} (Type 2): Stream redirection; produced by \textbf{\texttt{<}} (input redirect), \textbf{\texttt{>}} (output overwrite), and \textbf{\texttt{>>}} (output append, internal token \texttt{+}).
        \item \textbf{\texttt{PIPE}} (Type 3): Pipeline execution; produced by the pipe symbol \textbf{\texttt{|}}.
        \item \textbf{\texttt{LIST}} (Type 4): Sequential execution; produced by the semicolon \textbf{\texttt{;}}.
        \item \textbf{\texttt{BACK}} (Type 5): Asynchronous background execution; produced by the ampersand \textbf{\texttt{\&}}.
    \end{enumerate}

    \item \textbf{Where is cd handled, and why is it not run as a separate program?}\\[0.2em]
    \textbf{Where it is handled:} \texttt{cd} is handled directly inside the main loop of \texttt{main()} in \texttt{user/sh.c} (lines 166--170), preceding the call to \texttt{fork1()}:
\begin{lstlisting}[style=c]
if (cmd[0] == 'c' && cmd[1] == 'd' && cmd[2] == ' ') {
  // Chdir must be called by the parent, not the child.
  cmd[strlen(cmd) - 1] = 0; // chop \n
  if (chdir(cmd + 3) < 0)
    fprintf(2, "cannot cd %s\n", cmd + 3);
}
\end{lstlisting}
    \textbf{Why it cannot be an external program:} In UNIX process architecture, every process maintains its own current working directory in kernel memory (in xv6, stored in \texttt{struct proc -> cwd}). A child process created via \texttt{fork()} receives a duplicate copy of the parent's \texttt{cwd}. If \texttt{cd} were an independent external binary (e.g., \texttt{/bin/cd}), the shell would fork a child process, the child would call \texttt{chdir()} to modify its own directory, and then exit. The parent shell process would remain unaffected in its original directory. Therefore, changing the shell's working directory \textbf{must} be performed by the parent process itself via a built-in command.

    \item \textbf{Search sh.c for any code that handles * or \$. What does that explain about Part 2?}\\[0.2em]
    Inspecting \texttt{user/sh.c} for \texttt{*} and \texttt{\$} reveals that neither character appears anywhere in the tokenization table or parsing routines:
\begin{lstlisting}[style=c]
char whitespace[] = " \t\r\n\v";
char symbols[] = "<|>&;()";
\end{lstlisting}
    Because \texttt{*} and \texttt{\$} are neither whitespace nor recognized syntax symbols, \texttt{gettoken()} treats them as standard alphanumeric characters (type \texttt{'a'}), packaging them directly into \texttt{argv}.\\
    This confirms the empirical observations from Part 2 and Part 3: xv6 does not expand \texttt{ls *} because it lacks globbing logic, and it outputs \texttt{\$HOME} literally because it has no variable lookup table.
\end{enumerate}

% ==============================================================================
\section{Conclusion}

Through this laboratory work, I investigated operating system concepts from the user-kernel interface, contrasting the production-grade GNU \texttt{bash} shell with the minimalist, pedagogical xv6 shell.

By fixing and verifying the \texttt{sleep} utility in xv6, I observed how user-space libraries and applications interact with evolving kernel system calls across ABI revisions. The empirical side-by-side comparison between \texttt{bash} and xv6 revealed the core responsibilities of a Unix shell: coordinating processes via \texttt{fork()}, managing execution images via \texttt{exec()}, synchronizing lifecycles via \texttt{wait()}, and orchestrating stream redirection via \texttt{pipe()} and \texttt{dup()}. 

By developing and executing the \texttt{args} argument-inspection utility across both environments, I established a clear separation of concerns in Unix system architecture. The experiment proved that pathname wildcard expansion, environment variable substitution, and quote removal are purely user-space syntactic preprocessing duties handled by the shell, rather than services provided by the kernel or individual applications. Finally, dissecting \texttt{user/sh.c} demystified shell internals, demonstrating how an elegant 500-line C program leverages recursive descent parsing to construct an abstract syntax tree of commands that executes with remarkable efficiency.

\end{document}
